\begin{table}[t]
\centering
\small
\begin{tabular}{llrrrr}
\toprule
Run & target set & params & bpw & $\Delta$NLL & avg.\ per 100M \\
\midrule
\texttt{ternary\_task\_v4} & in\_proj\_qkv (18) & 113,246,208 & 1.7251 & $+0.088512$ & $0.078$ \\
\texttt{mlp\_task\_v1} & MLP (72) & 264,241,152 & 1.7250 & $+0.552558$ & $0.209$ \\
\texttt{mixed\_alloc\_v1} & both (90) & 377,487,360 & 1.7250 & $+0.637829$ & $0.169$ \\
\texttt{fullmodel\_v1} & all non-embedding (186) & 497,614,848 & 1.7250 & $+1.000022$ & $0.201$ \\
\bottomrule
\end{tabular}
\caption{Damage against coverage at ternary rate. Every row is asserted at generation time to share the same BF16 baseline to full double precision ($3.436710$) and the same target count (261,284); a matching count alone would not establish stream identity, but a loss computed over the whole stream agreeing in all digits does. The two disjoint families give an apparent additive reference, $0.088512 + 0.552558 = 0.641069$ against a measured $0.637829$, but these rows fit their codebooks on different pools; the frozen-artifact test reports an interaction of $+0.003818$ $[+0.000769,+0.006835]$, so damage is slightly super-additive. The jump at full coverage is composition, not breadth: going from the 90-tensor set to all 186 adds 120,127,488 parameters for $+0.362193$, i.e.\ $0.302$ per 100M, the highest marginal cost in the model. The last column is an average over each row's own target set, not a marginal rate.}
\label{tab:additivity}
\end{table}
